\section{The common domain and the exact chart integral}
\label{gc-domain}

Fix a compact connected semisimple matrix group $G$, a positive
invariant color inner product, and real orthonormal color coordinates.
The homogeneous current lattice has spacing one, prior classical
refinement is $M=L^k$, and the next blocking factor is fixed odd $L$.
Use the source geometry and ordinary averaged contours of Notebooks
19--22. All constants below may depend on $L$ and $G$, but not on $M$
or compatible periodic volume. The identities concern the established
small-potential domains, including their smaller complex continuations.
They do not define the classical effective action outside those domains.

Write $V=e^{ia}$ for a current background and $W=e^{iw}$ for next-coarse
data. The physical branch $V_*(W)$ satisfies
\begin{equation}\label{gc-physical-branch}
 M_L(V_*(W))=W,\qquad P_{V_*(W)}(y,x)=1.
\end{equation}
Notebook 22 identifies it with the unique small stationary solution
of the classical action plus the actual absolute averaged-contour
gauge penalty. Its strict support series is also available there.
The action $S_M$ is the ordinary classical constrained minimum of
the original Wilson action, with the coupling excluded.

For any sufficiently small independent background $V$, let
\begin{equation}\label{gc-constraint}
 F_V(B;c)=\frac1i\log\!\left[
 M_L(e^{iB}V)(c)M_L(V)(c)^{-1}\right].
\end{equation}
The exact nonlinear pivot chart $\Phi_V(z)$ is characterized by
\begin{equation}\label{gc-chart}
 F_V(\Phi_V(z))=0,\qquad
 \Phi_V(z)|_{\mathrm{retained}}=z.
\end{equation}
Each constraint reads its own pivot and no other pivot. This gives
uniqueness of its small solution with prescribed retained coordinates.
Denote its tangent insertion by $K_V=D\Phi_V(0)$ and its pivot
derivative at zero by
\begin{equation}\label{gc-pivot}
 A_{c,p}(V)=D_{p(c)}F_V(0;c),\qquad
 \mathcal D(V)=\prod_c\det A_{c,p}(V)^{-1}.
\end{equation}
The real determinants have the positive orientation fixed at the
identity background.

Let $\mathscr G$ be the absolute exponential gauge action and put
\begin{align}
 U_V(z)&=e^{i\Phi_V(z)}V, \notag\\
 T_V(z)&=S_M\!\left(\frac1i\log U_V(z)\right)
                         +\mathscr G(U_V(z)),             \label{gc-total-action}\\
 A_M(V)&=D_z^2T_V(0),\qquad A_{M,0}=A_M(1).              \label{gc-precision}
\end{align}
Thus $A_M$ includes both the constraint curvature and the logarithmic
coordinate curvature in the classical action. Away from the physical
branch the gauge Hessian is also its full second derivative. These
are the analytic centered product-domain extensions used in
Notebooks 20--21, not a Hessian obtained by dropping a background
current term.

Subtract the entire Taylor polynomial through degree two:
\begin{equation}\label{gc-taylor-remainder}
 \mathcal R_V(z)=T_V(z)-T_V(0)-DT_V(0)z
                           -\tfrac12\langle z,A_M(V)z\rangle.
\end{equation}
Let $j(X)$ be the Haar exponential-coordinate density divided by
its value at zero, and define the relative measure factor
\begin{equation}\label{gc-density}
 J_V(z)=
 \prod_b j(\Phi_V(z)(b))
 \prod_c\frac{\det A_{c,p}(V)}
 {\,\det D_{p(c)}F_V(\Phi_V(z);c)\,}.
\end{equation}
It is exactly the relative Haar density times the ambient
straightening determinant, as proved by pivot elimination in
Notebook 19. The rectangular derivative $D\Phi_V$ is not used as a
square Jacobian.

For positive real $g$, use the fixed retained window
\begin{equation}\label{gc-cutoff}
 \chi_R(\xi)=\prod_{b\ {\rm retained}}
                    \mathbf1_{\{|\xi(b)|_{\mathfrak g}\le R\}}.
\end{equation}
These are invariant Euclidean color balls at each retained bond,
not a product of intervals in individual color coordinates. A scalar
complex polydisc used for a Cauchy estimate is a different object.
The normalized integral is
\begin{equation}\label{gc-normalized-integral}
 Z_{M,g}(V)=\int d\mu_{A_M(V)^{-1}}(\xi)\,
      \chi_R(\xi)\,
      e^{-g^{-2}\mathcal R_V(g\xi)}J_V(g\xi).
\end{equation}
The real Gaussian is normalized. For complex backgrounds its density
is continued on the same real integration variables, with the
determinant branch fixed at the real positive reference. The bounds
of Notebook 21 make this analytic and nonvanishing on a common
small product domain, for its stated small-coupling regime.
At the physical branch $DT_V(0)=0$ and $\mathscr G(V)=0$; this
integral therefore supplies precisely the correction in the exact
small-field chart formula.

\section{Gauge covariance of the stationary background and pivot chart}
\label{gc-covariance}

For a next-coarse site gauge $u(y)$, its block-constant current lift is
\begin{equation}\label{gc-block-lift}
 u^\uparrow(x)=u(y),\qquad x\in B_L(y).
\end{equation}
Use the convention
$V^v(b)=v(b_-)V(b)v(b_+)^{-1}$ and the analogous convention for $W$.
We restrict $W$ and $u$ so all transformed fields, comparison branches,
and the gauge paths used below stay in one common small domain.
For example, the current bound
$\|\log V^{u^\uparrow}/i\|_\infty\le
C(\|a\|_\infty+\|\log u/i\|_\infty)$ and the prior uniform
branch bounds give a fixed smaller choice of radii.

\begin{lemma}[Block lift and stationary branch]\label{gc-branch-lemma}
For such real fields and gauges,
\begin{align}
 P_{V^{u^\uparrow}}(y,x)&=u(y)P_V(y,x)u(y)^{-1},\notag\\
 M_L(V^{u^\uparrow})&=M_L(V)^u,\qquad
 V_*(W^u)=V_*(W)^{u^\uparrow}.                           \label{gc-branch-equivariance}
\end{align}
These identities continue holomorphically in the corresponding
smaller complex parameter domain.
\end{lemma}
\begin{proof}
Every internal root contour lies in one block. Its endpoint
covariance is therefore conjugation by the same $u(y)$, and the
analytic group mean defining $P$ has that covariance. The ordinary
coarse average has endpoint covariance by its construction.

The local gauge identity for $S_M$ was proved in Notebook 16 by two
minimum comparisons, using a smooth original-fine extension of
arbitrary current-site gauge data. It applies to $u^\uparrow$:
the input is specified at current spacing one, and the smooth
extension has a norm bounded by the supremum of those site data
uniformly in $M$. A jump between current blocks does not produce
an $M$-dependent derivative in that extension.
Also $\mathscr G(V^{u^\uparrow})=\mathscr G(V)$, since the
contour matrices are conjugated and its defining trace is invariant.

Consequently the gauge diffeomorphism transports constrained
stationarity of $S_M+\mathscr G$ at $V_*(W)$ to constrained
stationarity at $V_*(W)^{u^\uparrow}$ with data $W^u$.
Its internal contours equal one by the first identity, and the
common-domain restriction places it in the same uniqueness
neighborhood as $V_*(W^u)$. The strictly positive retained Hessian
and stationary uniqueness established in Notebook 22 identify
the two configurations. No assertion that every stationary
configuration has $P=1$ is needed.
All the maps are holomorphic on the smaller complex domains.
Equality on an open real form extends coordinate by coordinate
to the connected complex parameter neighborhood.
\end{proof}

Put $V'=V^{u^\uparrow}$ and define the full relative bond action
and its retained restriction by
\begin{equation}\label{gc-adjoint-maps}
 (R_u B)(b)=\operatorname{Ad}_{u^\uparrow(b_-)}B(b),
 \qquad D_u=R_u|_{\mathrm{retained}}.
\end{equation}
For real compact gauges these maps are orthogonal in each color
fiber and have determinant one on their respective coordinate
spaces. Indeed an orthogonal adjoint matrix has determinant in
$\{1,-1\}$, and connectedness of $G$ fixes it to one.

\begin{lemma}[Exact nonlinear chart equivariance]\label{gc-chart-lemma}
On the common chart,
\begin{align}
 F_{V'}(R_uB;c)&=\operatorname{Ad}_{u(c_-)}F_V(B;c),       \label{gc-constraint-equivariance}\\
 \Phi_{V'}(D_uz)&=R_u\Phi_V(z),\qquad
 K_{V'}D_u=R_uK_V.                                      \label{gc-chart-equivariance}
\end{align}
\end{lemma}
\begin{proof}
One has $e^{iR_uB}V'=(e^{iB}V)^{u^\uparrow}$. Both ordinary
averages in the relative product defining $F$ transform with
the same endpoint gauges. Their product is conjugated at the
initial coarse endpoint. Conjugacy of the small logarithm proves
\eqref{gc-constraint-equivariance}.

The candidate $R_u\Phi_V(z)$ solves the transformed constraint
and has retained coordinates $D_uz$. The unique small pivot
solution is therefore $\Phi_{V'}(D_uz)$. Differentiating this
identity at zero gives the tangent identity.
\end{proof}

\section{Covariance of every normalization factor}
\label{gc-factors}

\begin{proposition}[Exact invariance of the normalized integral]
\label{gc-integral-proposition}
For real small $u$ and common backgrounds as above,
\begin{align}
 A_M(V')&=D_u A_M(V)D_u^T,&
 \mathcal D(V')&=\mathcal D(V),                           \label{gc-precision-delta}\\
 Z_{M,g}(V')&=Z_{M,g}(V).                                \label{gc-z-invariance}
\end{align}
Their analytic identities extend to the stated common complex domains.
\end{proposition}
\begin{proof}
Action gauge invariance, conjugacy of the absolute gauge action,
and Lemma~\ref{gc-chart-lemma} imply
\begin{equation}\label{gc-scalar-action-equivariance}
 T_{V'}(D_uz)=T_V(z).
\end{equation}
The constant, linear, and quadratic Taylor parts transform
separately under the fixed linear map $D_u$. Thus
\begin{equation}\label{gc-remainder-equivariance}
 \mathcal R_{V'}(D_uz)=\mathcal R_V(z).
\end{equation}
The quadratic identity reads $D_u^T A_M(V')D_u=A_M(V)$,
which gives the first formula in \eqref{gc-precision-delta}.
This argument also applies to the centered extension at independent
backgrounds that are not stationary.

Differentiating \eqref{gc-constraint-equivariance} in the pivot of
$c$ gives, at a general small $B$,
\begin{equation}\label{gc-pivot-equivariance}
 D_{p(c)}F_{V'}(R_uB;c)=
 \operatorname{Ad}_{u(c_-)}
 D_{p(c)}F_V(B;c)
 \operatorname{Ad}_{u^\uparrow(p(c)_-)}^{-1}.
\end{equation}
Both outside factors have determinant one. Evaluating at $B=0$
proves invariance of $\mathcal D(V)$. Evaluating at $\Phi_V(z)$
proves invariance of each normalized straightening determinant
in \eqref{gc-density}. Haar measure is invariant under conjugation,
and $\det\operatorname{Ad}_u=1$ in these real coordinates, so
$j(\operatorname{Ad}_uX)=j(X)$. Together with the chart identity
this proves
\begin{equation}\label{gc-measure-equivariance}
 J_{V'}(D_uz)=J_V(z).
\end{equation}

Make the real change of retained variables $\xi'=D_u\xi$.
It has unit Jacobian, carries the normalized Gaussian with
precision $A_M(V)$ to the one with precision $A_M(V')$, and
preserves every color norm in \eqref{gc-cutoff}. Equations
\eqref{gc-remainder-equivariance} and
\eqref{gc-measure-equivariance} prove \eqref{gc-z-invariance}.
This uses a real orthogonal change of variables; it neither
complexifies the sharp indicator nor deforms its integration domain.

Each scalar identity is holomorphic in the background and gauge
parameters on a smaller common domain. The real-form identity
theorem then supplies the complex identity. For complex adjoint
matrices the relevant algebraic transpose is bilinear; no claim
of unitary invariance of a complex sharp window is made.
\end{proof}

Define the separate Gaussian correction and total quantum correction by
\begin{align}
 \mathcal Q_M(V)&=-\tfrac12\log
                \det\!\left(A_M(V)A_{M,0}^{-1}\right),       \label{gc-q}\\
 \mathcal E_{M,g}^{\mathrm{eff}}(W)
   &=\log\frac{\mathcal D(V_*(W))}{\mathcal D(1)}
        +\mathcal Q_M(V_*(W))
        +\log Z_{M,g}(V_*(W)).                             \label{gc-effective}
\end{align}
All logarithms use compatible analytic branches. The determinant
branches are fixed by positivity and the identity reference;
the integral branch is the convergent connected expansion
from Notebook 21. On real sufficiently small fields that
integral is strictly positive, so this branch agrees with its
real logarithm.

\begin{theorem}[Local gauge invariance of the full correction]
\label{gc-effective-theorem}
In the common small domain,
\begin{equation}\label{gc-effective-invariance}
 \mathcal E_{M,g}^{\mathrm{eff}}(W^u)
                      =\mathcal E_{M,g}^{\mathrm{eff}}(W).
\end{equation}
The correction includes the pivot delta normalization and Gaussian
determinant, as well as the fluctuation integral. The explicit
field-independent prefactor and classical stationary action in the
full chart formula remain separate.
\end{theorem}
\begin{proof}
Lemma~\ref{gc-branch-lemma} identifies the two backgrounds.
Proposition~\ref{gc-integral-proposition} proves invariance of
$\mathcal D$, $\det A_M$, and $Z_{M,g}$ there.
The fixed scalar denominator $\det A_{M,0}$ is unchanged;
the matrix $A_{M,0}$ itself need not be invariant under a local
gauge. The compatible real logarithms consequently agree.
Holomorphic uniqueness gives the stated complex identity.
\end{proof}

Notebook 22 expresses the sum in \eqref{gc-effective} by strictly
supported coarse activities after reducing to a common positive
support weight and the stated smaller analytic radius. The theorem
proves invariance of that total sum. An arbitrary chosen individual
activity, or the deterministic connecting tree used to assign its
carrier, need not be gauge invariant.

\begin{proposition}[An inherited invariant functional]
\label{gc-inherited}
Suppose an inherited current functional $\mathcal E(U)$ is
holomorphic and gauge invariant on the needed common domain,
and insertion of $e^{\mathcal E(U_V(g\xi))}$ in
\eqref{gc-normalized-integral} gives a convergent integral
$Z_{M,g}^{\mathcal E}(V)$. Then
\begin{equation}\label{gc-inherited-integral}
 Z_{M,g}^{\mathcal E}(V^{u^\uparrow})
                          =Z_{M,g}^{\mathcal E}(V).
\end{equation}
If additionally this integral is nonvanishing and a compatible
holomorphic logarithm is chosen on the connected common domain,
the corresponding total logarithmic correction is locally
gauge invariant.
\end{proposition}
\begin{proof}
The chart identity gives
$U_{V'}(D_uz)=U_V(z)^{u^\uparrow}$, so the inherited factor is
unchanged in the same real change of variables. This proves
the integral identity. For the logarithm, the possible difference
is locally a constant in $2\pi i\mathbb Z$ and is zero at the
identity gauge; continuous paths in the stated common branch
therefore give equality. A verified convergent cluster expansion
is one way to supply the separate nonvanishing hypothesis.
\end{proof}
The additional hypothesis is material: convergence of a complex
integral alone does not exclude zeros. This proposition supplies
symmetry once the inherited integral is controlled. It does not
supply its required activity norm or repeated-step bounds.

\section{A finite-regulator Ward identity and its kernel consequence}
\label{gc-ward-section}

Let
\begin{equation}\label{gc-f}
 F(w)=\mathcal E_{M,g}^{\mathrm{eff}}(e^{iw}),
\end{equation}
or another holomorphic total coarse functional with the same
local gauge invariance. Here $w$ is the next-coarse link angle
in next-coarse unit coordinates.

\begin{theorem}[Quadratic Ward identity]\label{gc-ward-theorem}
On every admissible finite periodic regulator,
\begin{equation}\label{gc-no-linear}
 DF(0)=0.
\end{equation}
For the signed unit coarse gradient
$(\partial\lambda)(b)=\lambda(b_-)-\lambda(b_+)$, one has
\begin{equation}\label{gc-ward}
 D^2F(0)[v,\partial\lambda]=0
 \quad\hbox{for every }v,\lambda,
 \qquad H_F\partial=0,\quad H_F=D^2F(0).
\end{equation}
\end{theorem}
\begin{proof}
Constant site gauges act by adjoint conjugation on each bond
angle. Differentiating their invariance at $w=0$ says that each
individual linear color covector is adjoint invariant.
Infinitesimally such a covector annihilates every commutator.
Since $\mathfrak g=[\mathfrak g,\mathfrak g]$ for a semisimple
Lie algebra, it is zero. Independent variations at each bond
prove \eqref{gc-no-linear}. Only constant gauges in a neighborhood
of the identity are needed for this reasoning.

For a site field $\lambda$ set $u_t=e^{it\lambda}$ and denote
the infinitesimal gauge velocity of the potential by
$X_\lambda(w)=\left.\frac{d}{dt}w^{u_t}\right|_{t=0}$.
Local multiplication gives $X_\lambda(0)=\partial\lambda$.
Gauge invariance states $DF(w)X_\lambda(w)=0$.
Differentiating in the direction $v$ at zero yields
\[
 D^2F(0)[v,\partial\lambda]
          +DF(0)\bigl(DX_\lambda(0)v\bigr)=0.
\]
The second term vanishes by \eqref{gc-no-linear}, proving
\eqref{gc-ward}. This argument concerns the total functional,
not the separate members of a support decomposition.
\end{proof}

For completeness, the transported activity norm gives an actual
exponentially summable Hessian. It is not necessary to assume this
kernel bound independently at a finite regulator.

\begin{lemma}[Hessian from connected analytic activities]
\label{gc-hessian-bound}
Suppose $F(w)=\mathrm{Vac}+\sum_K F_K(w|_K)$ with nonempty connected
carriers, $d_0$ scalar coordinates per unit cell, holomorphy on a
scalar polydisc of radius $r$, and
\begin{equation}\label{gc-activity-hypothesis}
 \sup_i\sum_{K\ni i} e^{a|K|}\|F_K\|_r\le\epsilon,\qquad a>0.
\end{equation}
For any $0\le\mu<a$, in the maximum-coordinate block operator norm,
\begin{equation}\label{gc-hessian-estimate}
 \sup_i\sum_j e^{\mu d(i,j)}\|(H_F)_{ij}\|
 \le \frac{2d_0\epsilon}{r^2}\,
          C_{a-\mu},\qquad
 C_\rho=\sup_{k\ge1}k e^{-\rho k}\le\frac1{e\rho}.
\end{equation}
An equivalent Euclidean block norm changes only the fixed
dimension constant.
\end{lemma}
\begin{proof}
Coordinatewise Cauchy bounds a second scalar partial derivative
of $F_K$ at zero by $2r^{-2}\|F_K\|_r$; the factor two covers
repeated differentiation in the same coordinate. Its $(i,j)$
block has norm at most $2d_0r^{-2}\|F_K\|_r$ and vanishes unless
$i,j\in K$. Connectedness gives $d(i,j)\le |K|-1$.
Summing its at most $|K|$ column cells and using
$|K|e^{\mu(|K|-1)}\le C_{a-\mu}e^{a|K|}$
proves \eqref{gc-hessian-estimate}. The positive bound justifies
the derivative and support sums.
\end{proof}

The bounds of Notebook 22 supply \eqref{gc-activity-hypothesis}
for the total correction on its smaller analytic domain.
This gives a uniform Hessian bound at fixed $g,L,r$ and the
allowed classical depth. The factor $r^{-2}$ is explicit;
this is not a coupling-flow estimate.

\begin{proposition}[Zero-momentum consequence under a kernel limit]
\label{gc-zero-momentum}
Assume additionally that a sequence of these Hessians is translation
invariant and converges, on tori exhausting all coordinate directions,
to a convolution kernel on $\mathbb Z^4$, with a uniform positive
exponential summability bound. Then its Fourier multiplier
$\widehat H_F(p)$ is analytic near real momentum zero and
\begin{equation}\label{gc-zero}
 \widehat H_F(p)d(p)=0,\qquad
 d_\nu(p)=1-e^{ip_\nu},\qquad \widehat H_F(0)=0.
\end{equation}
The first equation includes the identity on color space.
\end{proposition}
\begin{proof}
Exponential summability gives locally uniform convergence of the
Fourier series in a complex strip. The finite-volume Ward identity
passes to the limiting kernel; since the gradient is a finite
difference operator, this also follows directly by taking limits
in its kernel equation. Fourier transformation gives the first
identity for continuous real $p$.
For $p=t e_\nu$, with $t\ne0$ small, the only nonzero component
of $d(p)$ is $1-e^{it}$ in direction $\nu$. Thus
$\widehat H_F(t e_\nu)e_\nu=0$ on every color vector.
Continuity as $t\to0$ annihilates every column direction of
$\widehat H_F(0)$, proving the last identity.
\end{proof}
The conclusion excludes a nonzero momentum-independent quadratic
mass coefficient in that limit. Finite discrete momenta alone would
not justify the continuous-momentum step. The additional symmetry
and kernel-limit hypotheses have therefore been retained explicitly.

\section{Use in the existence argument}
\label{gc-next}

The complete defined small-field step preserves local gauge
invariance, including its cutoff and normalization factors.
Together with the strict coarse support construction, this gives
an analytic local correction whose total quadratic kernel obeys
the Ward identity. The inherited-functional proposition shows
how the same symmetry can pass through a later integration when
its analytic norm and nonvanishing are established.

The next task is to extract the vacuum and marginal gauge-action
coefficient and control the remaining interactions under repeated
scaling. Transversality does not determine that coefficient or its
flow. Neither does it alone classify all allowed momentum terms:
the rotation and reflection properties used in such a classification
must be proved for the actual construction. Stable analytic or
gauge-orbit domains, inherited fluctuation-interaction estimates,
large-field complements, and observable insertion bounds remain
part of the full existence proof. This notebook proves the symmetry
input to that analysis.
